% =====================================================================
% 01_hypothesis.tex  --  Hypothesis (Exp. 4: d-analysis + noise proxy)
% =====================================================================

\section{Hypothesis}
\label{sec:dproxy-hypothesis}

Contrastive decoding, in the visual forms VCD~\citep{leng2024vcd} and
SID~\citep{huo2024sid}, builds its output score as
$C = (1+\alpha)E - \alpha A$, which at $\alpha = 1$ equals $2E - A = E + d$,
where $d = E - A$ is the gap between the clean-image expert logit $E$ and the
degraded-image amateur logit $A$. The quantity $d$ is therefore exactly the
adjustment that contrastive decoding adds to the plain expert score at every
candidate token. The stated purpose of the method is to push down tokens that
the language prior favours without visual support, so for a hallucinated object
$d$ should be negative.

We test this premise in two complementary ways.

\begin{description}
  \item[\textbf{H1 (Selective suppression).}] If contrastive decoding suppresses
  hallucinations, the adjustment $d$ should be more negative for hallucinated
  objects than for real objects. We measure the sign and size of $d$ separately
  for the two groups. A working suppressor predicts clearly negative $d$ on
  hallucinations, and more negative on hallucinations than on real objects.
  \item[\textbf{H2 (Mechanism versus perturbation).}] If the change contrastive
  decoding produces in the hallucination rate is a targeted effect of the
  amateur branch, then replacing that branch with random noise of the same
  magnitude should not reproduce it. We construct noise proxies whose only input
  is expert logits plus a random bonus matched to the measured size of $d$, and
  compare their CHAIR scores against the real methods.
\end{description}

\noindent The first test asks whether the contrastive adjustment points in the
direction the method claims. The second asks whether the measured effect on
hallucination requires the amateur branch at all, or whether generic noise of
the same magnitude is enough to explain it. We run both on two model families
(LLaVA-1.5-7B~\citep{liu2024llava} and
Qwen2.5-VL-7B~\citep{bai2025qwen25vl}) and both contrastive methods (VCD and
SID). This complements the benchmark-level critique of
\citet{yin2025mirage} by examining the logit adjustment itself.
